This paper presents an evidence-grounded, simulation-oriented planning pipeline built around three implemented components: Memento-style case retrieval, a Hope adaptive weighting controller, and a reflection loop for iterative prompt refinement. The system generates structured joint-operation plans, evaluates them through a five-stage kill-chain simulation model, and exports the selected plan into Department of Defense Architecture Framework (DoDAF) OV-5b, DoDAF OV-6c, and Command and Control Systems -- Simulation Systems Interoperation (C2SIM) artifacts. To avoid unsupported claims, this study only describes mechanisms that are explicitly present in the implemented system and experiment outputs, and frames itself as an application-oriented integration paper rather than as a claim of a new standalone planning theory.

The main single-scenario ablation study uses a sample coastal joint-assault scenario with 10 optimization iterations, 50 Monte Carlo simulation runs, a fixed random seed of 7, and writeback disabled. Under the default HOPE parameters ($\theta=0.5$, $\epsilon=0.02$), the full pipeline improves mission success from 62.75 to 65.76, overall effectiveness from 73.84 to 78.75, and the loss-exchange ratio from 1.71 to 2.19 relative to the pure large language model baseline on an NVIDIA GeForce RTX 4090 D GPU. A multi-seed replication across five seeds confirms the Full pipeline's gain is statistically significant: the paired mean mission-success gain is +3.81 points ($\text{SD}=2.27$, paired $t(4)=3.75$, $p=0.02$, Cohen's $d=1.68$), and the overall-effectiveness gain reaches $p<0.01$ ($t(4)=5.87$). The cross-scenario transfer study uses five scenario families, 20 planning iterations, and leave-one-source-scenario-out case banks derived from 250 source cases. In that evidence set, the full pipeline outperforms the pure baseline in four of five scenarios, increasing average mission success from 65.01 to 66.48 and average overall effectiveness from 72.49 to 77.92, while the relative ordering among the Memento-only and Hope-only variants remains scenario-dependent. A separate 50-iteration suite is treated as an auxiliary convergence view rather than mixed into the main transfer table.

The paper also documents the remaining validity limits, including scenario dependence in module gains and the need for per-component multi-seed extension beyond the current single-scenario-family evidence. A three-backend comparison and a multi-seed Reflection-only replication (3 seeds) confirm that the pipeline's gains are not artifacts of a single model or seed choice. In particular, the reported confidence intervals reflect within-seed Monte Carlo variance under the fixed seed rather than across-seed variance. The refreshed canonical result bundles now serialize a top-level runtime manifest that records the local backend identifier, software stack, available CPU/GPU metadata, and wall-clock timing summaries for the reported evidence.
